\documentclass[11pt,a4paper]{article}

% ─────────────────────────────────────────────────────────────────────────────
%  Exp Intelligence — Technical & Strategic Report
% ─────────────────────────────────────────────────────────────────────────────

\usepackage[utf8]{inputenc}
\usepackage[T1]{fontenc}
\usepackage[margin=1in]{geometry}
\usepackage{lmodern}
\usepackage{microtype}
\usepackage{graphicx}
\usepackage{xcolor}
\usepackage{titlesec}
\usepackage{enumitem}
\usepackage{booktabs}
\usepackage{longtable}
\usepackage{array}
\usepackage{tabularx}
\usepackage{fancyhdr}
\usepackage{listings}
\usepackage{tcolorbox}
\usepackage{tikz}
\usetikzlibrary{arrows.meta, positioning, shapes.geometric, calc}
\usepackage{hyperref}

% ── Colours ──────────────────────────────────────────────────────────────────
\definecolor{brandblue}{HTML}{1F4E79}
\definecolor{brandteal}{HTML}{2E8B8B}
\definecolor{lightgray}{HTML}{F4F6F8}
\definecolor{codebg}{HTML}{F7F7F9}
\definecolor{rulegray}{HTML}{C9D2DA}
\definecolor{accent}{HTML}{C0563B}

\hypersetup{
  colorlinks=true,
  linkcolor=brandblue,
  urlcolor=brandteal,
  citecolor=brandblue,
  pdftitle={Exp Intelligence — Technical and Strategic Report},
  pdfauthor={Aadi Jain}
}

% ── Section styling ──────────────────────────────────────────────────────────
\titleformat{\section}
  {\normalfont\Large\bfseries\color{brandblue}}
  {\thesection}{0.6em}{}
\titleformat{\subsection}
  {\normalfont\large\bfseries\color{brandteal}}
  {\thesubsection}{0.6em}{}
\titleformat{\subsubsection}
  {\normalfont\normalsize\bfseries\color{black}}
  {\thesubsubsection}{0.6em}{}

% ── Headers / footers ────────────────────────────────────────────────────────
\pagestyle{fancy}
\fancyhf{}
\renewcommand{\headrulewidth}{0.4pt}
\renewcommand{\footrulewidth}{0pt}
\fancyhead[L]{\small\color{brandblue}Exp Intelligence}
\fancyhead[R]{\small\color{brandblue}Technical \& Strategic Report}
\fancyfoot[C]{\thepage}

% ── Code listing style ───────────────────────────────────────────────────────
\lstset{
  basicstyle=\ttfamily\footnotesize,
  backgroundcolor=\color{codebg},
  frame=single,
  rulecolor=\color{rulegray},
  breaklines=true,
  breakatwhitespace=true,
  showstringspaces=false,
  keywordstyle=\color{brandblue}\bfseries,
  commentstyle=\color{brandteal}\itshape,
  stringstyle=\color{accent},
  numbers=none,
  xleftmargin=6pt,
  xrightmargin=6pt,
  aboveskip=8pt,
  belowskip=8pt,
  columns=fullflexible,
}

% ── Callout box ──────────────────────────────────────────────────────────────
\newtcolorbox{keybox}[1]{
  colback=lightgray,
  colframe=brandteal,
  boxrule=0.6pt,
  arc=2pt,
  left=8pt,right=8pt,top=6pt,bottom=6pt,
  title=#1,
  fonttitle=\bfseries\color{white},
  coltitle=white,
  colbacktitle=brandteal
}

\newcolumntype{Y}{>{\raggedright\arraybackslash}X}

\setlength{\parskip}{0.5em}
\setlength{\parindent}{0pt}

% ─────────────────────────────────────────────────────────────────────────────
\begin{document}

% ── Title page ───────────────────────────────────────────────────────────────
\begin{titlepage}
\centering
\vspace*{1.4cm}
{\large\color{brandblue}\bfseries GlobalLogic \textnormal{(A Hitachi Group Company)}}\\[0.15cm]
{\small\color{brandteal} BSC-NGT Portfolio --- GenAI Consultancy}\\[0.8cm]
{\color{brandblue}\rule{\textwidth}{2pt}}\\[0.6cm]
{\Huge\bfseries\color{brandblue} Exp Intelligence}\\[0.4cm]
{\LARGE AI-Assisted Storefront Optimization Pipeline}\\[0.3cm]
{\large\itshape Technical Architecture, Rationale \& Strategic Roadmap}\\[0.4cm]
{\color{brandblue}\rule{\textwidth}{2pt}}\\[0.8cm]
{\large\bfseries by Aadi Jain}\\[1.0cm]

\begin{minipage}{0.85\textwidth}
\centering\itshape\large
``Turn raw storefront analytics into ranked, data-grounded improvement
suggestions --- each with a ready-to-paste code patch --- reviewed by a
product owner and picked up by developers.''
\end{minipage}

\vfill

\begin{tabular}{rl}
\textbf{Author:} & Aadi Jain --- Backend \& Backlog\\
\textbf{Programme:} & GenAI Consultancy Internship, BSC-NGT\\
\textbf{Document type:} & Technical \& strategic report\\
\textbf{Scope:} & Context, Why, How, PM \& delivery, Improvements, Future, Deployment\\
\textbf{Services covered:} & Shopify Extractor, Analytics, AI, Review Hub\\
\textbf{Date:} & July 2026\\
\end{tabular}

\vspace{1cm}
{\color{rulegray}\rule{\textwidth}{0.5pt}}
\end{titlepage}

% ── TOC ──────────────────────────────────────────────────────────────────────
\tableofcontents
\thispagestyle{fancy}
\newpage

% ═════════════════════════════════════════════════════════════════════════════
\section{Executive Summary}
% ═════════════════════════════════════════════════════════════════════════════

\textbf{Exp Intelligence} is a microservice pipeline that closes the loop
between \emph{observing} how a Shopify storefront performs and \emph{acting}
on it. It automatically pulls a store's own analytics, detects statistically
significant problems, uses large language models (LLMs) to translate those
problems into prioritized, plain-English improvement suggestions --- each
carrying a concrete code patch --- and presents them in a review dashboard
where a Product Owner (PO) approves or rejects them and developers pick up the
approved ones.

The system is composed of four independently deployable services:

\begin{center}
\begin{tabularx}{\textwidth}{@{}l l c Y@{}}
\toprule
\textbf{Service} & \textbf{Folder} & \textbf{Port} & \textbf{Responsibility}\\
\midrule
Shopify Extractor & \texttt{shopify-pp} & 3000 & OAuth install on a store; extract ShopifyQL metrics + real theme code; forward to Analytics\\
Analytics Service & \texttt{analytic service} & 4000 & Normalize \& ingest data; run a Playwright crawler; detect problems statistically\\
AI Service & \texttt{ai service} & 5001 & Turn analytics + screenshots into ranked suggestions and code patches via LLMs\\
Review Hub & \texttt{review-hub} & 5173 & React dashboard: PO approve/reject, developer code view\\
\bottomrule
\end{tabularx}
\end{center}

Data flows strictly in one direction:
\[
\texttt{shopify-pp}\;\rightarrow\;\texttt{analytic service}\;\rightarrow\;\texttt{ai service}\;\rightarrow\;\texttt{review-hub}
\]

\begin{keybox}{The core value proposition}
Store owners drown in dashboards but starve for decisions. Exp Intelligence
does not add another dashboard --- it produces \emph{ranked, data-grounded
actions with the code to implement them}, gated by a human review step so that
nothing ships without a person approving it.
\end{keybox}

% ═════════════════════════════════════════════════════════════════════════════
\section{Project Context \& Organization}
% ═════════════════════════════════════════════════════════════════════════════

Exp Intelligence (also referred to as the \emph{Experience Intelligence /
Analytics-to-UX Recommendation Engine}) was built during a GenAI Consultancy
internship engagement within GlobalLogic's \textbf{BSC-NGT} (Business Services
\& Consulting --- Next Generation Technologies) portfolio. It targets D2C
(direct-to-consumer) Shopify merchants in the fashion, beauty, and lifestyle
segment.

\subsection{Team \& stakeholders}

\begin{center}
\begin{tabularx}{\textwidth}{@{}l l Y@{}}
\toprule
\textbf{Name} & \textbf{Role} & \textbf{Relevance to the project}\\
\midrule
Ambuj Prakash & Project Mentor \& Lead & Reviews deliverables, overall engagement ownership\\
Nikhil Gupta & Technical Mentor & Technical architecture \& implementation guidance\\
Shubham Tomar & Mentor & Onboarding, guidance, reviews\\
Akanksha Rajwanshi & Mentor & Onboarding, guidance, reviews\\
Samad Khan & Team / BSC-NGT & BSC-NGT Portfolio team\\
Mansimar & Peer --- Backend \& Backlog & Co-owned backend (Analytics + AI services) \& backlog\\
\textbf{Aadi Jain} & Peer --- Backend \& Backlog & Co-owned backend (Analytics + AI services) \& backlog\\
Pranjal Tanwar & Peer --- Frontend \& Backlog & Co-owned frontend (Admin Portal / Review Hub) \& backlog\\
Rishabh Varshney & Peer --- Frontend \& Backlog & Co-owned frontend (Admin Portal / Review Hub) \& backlog\\
\bottomrule
\end{tabularx}
\end{center}

\subsection{The v1 success criterion}

The success criterion agreed with the team is deliberately narrow and
outcome-based: \textbf{one real merchant pilot where at least one AI-generated
recommendation is actually implemented on the merchant's store and its impact
measured.} Everything in this report --- the architecture, the review gate, the
code generation --- exists to make that single outcome achievable and
repeatable.

\subsection{Source repositories}

Each service lives in its own GitHub repository (assembled as submodules of the
\texttt{exp-intelligence} monorepo):

\begin{center}
\begin{tabularx}{\textwidth}{@{}l Y@{}}
\toprule
\textbf{Service} & \textbf{Repository}\\
\midrule
Shopify Extractor & \url{https://github.com/aadij1905/shopify-extractor}\\
Analytics Service & \url{https://github.com/aadij1905/analytics-service}\\
AI Service & \url{https://github.com/aadij1905/ai-service}\\
Review Hub & \url{https://github.com/aadij1905/review_hub}\\
\bottomrule
\end{tabularx}
\end{center}

\begin{keybox}{A note on this document vs.\ the handover snapshot}
Parts of this report draw on the internship \emph{handover} document. The
\emph{technical} sections here reflect the \textbf{current codebase}: all four
backlog epics --- including E-004 code generation --- are complete and in MVP
scope, and suggestion generation supports \emph{multiple} LLM providers with a
generate$\rightarrow$judge orchestration rather than a single fixed model.
Where the handover and the code differ in incidental detail (service ports, the
``Admin Portal''\,/\,``Review Hub'' name, the default model), this report
follows the code.
\end{keybox}

% ═════════════════════════════════════════════════════════════════════════════
\section{Why --- The Problem \& Motivation}
% ═════════════════════════════════════════════════════════════════════════════

\subsection{The two-buyer problem}

The engagement was framed around a \textbf{two-buyer problem} that shaped most
product decisions:

\begin{itemize}[leftmargin=1.4em]
  \item \textbf{The Growth Lead / marketer} feels the conversion (CRO) pain
  daily but has \emph{no access to the store's codebase} to act on it.
  \item \textbf{The Developer} controls the codebase and can make changes, but
  has \emph{no visibility} into analytics or CRO priorities.
\end{itemize}

Neither can close the loop alone. This is why the product deliberately runs two
tracks --- an insight/prioritization/business-impact experience for the Growth
Lead (the PO view), and an implementation-ready experience with code patches
for the Developer --- and why a human approval step sits between them. The
sections below expand the underlying analytics and engineering rationale.

\subsection{The problem}

Shopify merchants have access to an enormous amount of analytics --- sessions,
conversion rates, bounce rates, traffic sources, device breakdowns, checkout
funnels, Core Web Vitals --- but this data is:

\begin{itemize}[leftmargin=1.4em]
  \item \textbf{Descriptive, not prescriptive.} A dashboard tells you the
  product page bounce rate is 71\%. It does not tell you \emph{why}, whether
  that is statistically meaningful given the traffic, or \emph{what to change
  in the code} to fix it.
  \item \textbf{Disconnected from implementation.} Even when a merchant knows
  \emph{what} is wrong, translating ``improve the product page CTA'' into an
  actual edit in a Liquid theme file is a developer task that rarely happens.
  \item \textbf{Noisy.} A wild-looking metric over six sessions is almost
  certainly noise, yet naive tools flag it anyway, eroding trust.
\end{itemize}

\subsection{The insight}

The gap is not \emph{more measurement} --- it is the translation layer between
\emph{measurement} and \emph{implementation}. Three capabilities, chained
together, close that gap:

\begin{enumerate}[leftmargin=1.4em]
  \item \textbf{Statistical rigor} to separate real problems from sampling
  noise (a one-sample $z$-test on proportions, not raw multipliers).
  \item \textbf{Multimodal reasoning} --- combining the \emph{numbers}
  (ShopifyQL) with what a page \emph{looks like} (crawler screenshots analyzed
  by a vision model) --- because many conversion problems (a CTA below the
  fold, missing social proof) are invisible in metrics alone.
  \item \textbf{Code generation grounded in the store's own theme}, so a
  suggestion arrives as a real patch against the real installed theme, not
  generic advice.
\end{enumerate}

\subsection{Why a human-in-the-loop design}

Every suggestion passes through a Product Owner before a developer sees it.
This is deliberate: LLM output is probabilistic, and shipping automated code
changes to a live storefront without review is unacceptable risk. The PO layer
turns the AI from an ``autonomous agent'' into a \emph{trusted assistant} ---
the AI proposes, a human disposes, and only then does a developer implement.
The roadmap (Section~\ref{sec:future}) describes how this gate is preserved
even as the system moves toward auto-applying patches.

% ═════════════════════════════════════════════════════════════════════════════
\section{How --- System Architecture}
% ═════════════════════════════════════════════════════════════════════════════

\subsection{Architecture at a glance}

\begin{figure}[h]
\centering
\resizebox{\textwidth}{!}{%
\begin{tikzpicture}[
  svc/.style={rectangle, rounded corners=3pt, draw=brandblue, line width=1pt,
    fill=lightgray, text width=3.1cm, align=center, minimum height=1.9cm},
  store/.style={cylinder, shape border rotate=90, aspect=0.28, draw=brandteal,
    line width=1pt, fill=white, text width=1.9cm, align=center, minimum height=1.6cm},
  aux/.style={rectangle, rounded corners=2pt, draw=rulegray, dashed, line width=0.8pt,
    fill=white, text width=3.0cm, align=center, font=\footnotesize},
  persona/.style={rectangle, rounded corners=2pt, draw=accent, line width=1pt,
    fill=white, text width=2.4cm, align=center, font=\footnotesize, minimum height=1cm},
  flow/.style={-{Latex[length=2.6mm]}, line width=1pt, brandblue},
  auxflow/.style={-{Latex[length=2.2mm]}, dashed, line width=0.8pt, gray},
  lbl/.style={midway, above, font=\footnotesize\itshape, text=black},
]

\node[store] (store) {\textbf{Shopify}\\ \textbf{Store}};
\node[svc, right=1.3cm of store] (ext) {\textbf{Shopify Extractor}\\[2pt] {\scriptsize shopify-pp \textbullet\ :3000}\\[2pt] {\scriptsize OAuth \textbullet\ ShopifyQL \textbullet\ theme code}};
\node[svc, right=1.3cm of ext] (an) {\textbf{Analytics Service}\\[2pt] {\scriptsize :4000}\\[2pt] {\scriptsize normalize \textbullet\ crawler \textbullet\ detectors}};
\node[svc, right=1.3cm of an] (ai) {\textbf{AI Service}\\[2pt] {\scriptsize :5001}\\[2pt] {\scriptsize suggestions \textbullet\ code patches}};
\node[svc, right=1.3cm of ai] (hub) {\textbf{Review Hub}\\[2pt] {\scriptsize :5173}\\[2pt] {\scriptsize React dashboard}};

\node[aux, below=1.0cm of an] (crawl) {Playwright crawler\\ desktop + mobile screenshots};
\node[aux, below=1.0cm of ai] (llm) {LLM providers\\ Groq \textbullet\ Gemini \textbullet\ Mistral \textbullet\ Cerebras};
\node[persona] (po) at ([yshift=0.75cm]$(hub.east)+(2.2cm,0)$) {Product Owner\\ approve / reject};
\node[persona] (dev) at ([yshift=-0.75cm]$(hub.east)+(2.2cm,0)$) {Developer\\ code patch};

\draw[flow] (store) -- (ext) node[lbl]{extract};
\draw[flow] (ext) -- (an) node[lbl]{ingest};
\draw[flow] (an) -- (ai) node[lbl]{data + flags};
\draw[flow] (ai) -- (hub) node[lbl]{report};
\draw[flow] (hub.east) -- (po.west);
\draw[flow] (hub.east) -- (dev.west);

\draw[auxflow] (crawl) -- (an);
\draw[auxflow] (llm) -- (ai);
\draw[auxflow] (ai.north) to[out=90, in=90, looseness=1.45]
  node[pos=0.5, above, font=\scriptsize\itshape, text=gray, yshift=1pt]{theme code (/code/generate)}
  (ext.north);

\end{tikzpicture}%
}
\caption{The Exp Intelligence pipeline. Solid arrows are the one-directional
data flow; dashed arrows are supporting dependencies (the crawler and LLM
providers) and the on-demand theme-code fetch used when generating a patch.}
\label{fig:arch}
\end{figure}

\subsection{End-to-end data flow}

\begin{enumerate}[leftmargin=1.4em]
  \item \textbf{Install \& extract} (\texttt{shopify-pp}). A merchant installs
  the embedded app on their store via Shopify OAuth. The extractor runs 10
  ShopifyQL queries (overview, sales, per-page, traffic sources, devices,
  funnel, checkout conversion, and three Core Web Vitals: LCP, CLS, INP) and
  can also pull the \emph{real theme code} that renders any given page.
  \item \textbf{Ingest \& normalize} (\texttt{analytic service}). Raw query
  rows are normalized into a uniform shape, then a headless Playwright crawler
  visits the storefront, takes desktop + mobile screenshots, and enriches page
  records (e.g.\ ``is the CTA above the fold?'', ``is there social proof?'').
  \item \textbf{Detect} (\texttt{analytic service}). Deterministic detectors
  run a $z$-test against the site baseline and emit ranked, severity-scored
  flags (high-bounce page, low-conversion page, poor traffic source, cart
  abandonment, funnel drop-off, revenue decline, poor Core Web Vitals, CTA
  below fold).
  \item \textbf{Generate} (\texttt{ai service}). The AI service pulls the
  normalized data + flags, and --- if screenshots exist --- feeds them to a
  vision model in one integrated pass. LLMs draft suggestions; a deterministic
  validator grounds every cited number; a revenue forecaster attaches a
  data-derived impact estimate.
  \item \textbf{Review} (\texttt{review-hub}). The PO sees suggestions
  \emph{without} code (problem, plain-English recommendation, effort, impact,
  confidence) and approves or rejects. Developers see only approved
  suggestions, each with a ready-to-paste code patch generated on demand
  against the store's real theme.
\end{enumerate}

\subsection{Service 1 --- Shopify Extractor (\texttt{shopify-pp})}

\textbf{Stack:} Node.js + Express, \texttt{better-sqlite3} for persistence.

\begin{itemize}[leftmargin=1.4em]
  \item \textbf{OAuth \& security.} Implements the full Shopify OAuth handshake
  with HMAC verification on both the install redirect and webhooks (raw-body
  HMAC, mounted before \texttt{express.json()} so signatures verify against the
  exact bytes). Access tokens are encrypted at rest in SQLite using a 32-byte
  \texttt{ENCRYPTION\_KEY}.
  \item \textbf{Compliance.} Handles Shopify's mandatory GDPR webhooks
  (\texttt{customers/data\_request}, \texttt{customers/redact},
  \texttt{shop/redact}). The app stores no customer PII --- only shop domain +
  access token --- so redaction simply deletes the shop record.
  \item \textbf{ShopifyQL extraction.} A single \texttt{SESSION\_FILTER}
  constant guarantees every query (\texttt{overview}, \texttt{pages},
  \texttt{traffic}, \texttt{devices}, \texttt{funnel}, \texttt{checkout})
  counts the \emph{same} session universe, so the cards never disagree.
  \item \textbf{Theme code extraction.} \texttt{/api/debug/theme-code} returns
  the actual Liquid files that render a given page --- the raw material the AI
  service uses to write store-specific patches.
\end{itemize}

\subsection{Service 2 --- Analytics Service (\texttt{analytic service})}

\textbf{Stack:} Node.js + Express, Playwright (Chromium), optional AWS S3 for
screenshot storage.

\begin{itemize}[leftmargin=1.4em]
  \item \textbf{Normalization.} Converts raw ShopifyQL rows into a stable
  internal shape; Core Web Vitals are bucketed against Google's own thresholds
  (LCP good $\le 2500$ms, CLS good $\le 0.1$, INP good $\le 200$ms).
  \item \textbf{Crawler.} A reusable headless Chromium session visits each
  significant page at desktop (1280$\times$800) and mobile (390$\times$844)
  viewports. It detects the add-to-cart CTA using known theme selectors
  \emph{plus} a theme-agnostic text heuristic (``add to cart / buy now / \dots''),
  measures whether it is above the fold, and detects social-proof widgets.
  \item \textbf{Statistical detectors.} The heart of the noise-rejection
  strategy. See the callout below.
\end{itemize}

\begin{keybox}{Why a $z$-test instead of a raw multiplier}
A page bouncing at 2$\times$ the site average over 6 sessions is probably
noise; the same rate over 6{,}000 sessions is a real signal. A fixed
``minimum sessions'' floor cannot tell these apart across store sizes. The
one-sample $z$-test for a proportion measures how many standard errors an
observed rate sits from the baseline:
\[
z = \frac{\hat{p} - p_0}{\sqrt{p_0(1-p_0)/n}}
\]
A flag only fires when \emph{both} the effect size (e.g.\ bounce $>1.3\times$
average) \emph{and} significance ($|z| \ge 1.645$, $\approx$95\% confidence)
hold. Flags on the same page are annotated as co-occurring so the AI treats
them as one causal story, not $N$ independent findings.
\end{keybox}

\subsection{Service 3 --- AI Service (\texttt{ai service})}

\textbf{Stack:} Node.js + Express, multi-provider LLM layer (Groq, Gemini,
Mistral, Cerebras, Ollama).

\begin{itemize}[leftmargin=1.4em]
  \item \textbf{Provider-agnostic core.} A single \texttt{llm.js} entry point
  resolves the provider from the \texttt{AI\_PROVIDER} env var and lazily loads
  only the SDK actually selected, so a missing optional dependency never breaks
  the service. All clients share a uniform \texttt{\{ content, meta \}} shape.
  \item \textbf{Generate$\rightarrow$Judge orchestration.} When
  \texttt{AI\_PROVIDER=orchestrate}, several generator models draft candidates
  in parallel; candidates are deduplicated by fingerprint
  (\texttt{category::page}, tracking cross-model agreement); a judge model
  merges, filters, and re-ranks them. Any generator that is rate-limited is
  simply skipped (\texttt{Promise.allSettled}); if the judge fails, the deduped
  candidates are returned so a report always comes back.
  \item \textbf{Vision-integrated path.} When the crawler produced
  screenshots, a vision model (Gemini) reasons over the images directly in the
  same pass that writes suggestions --- a visual-flaw pre-pass also injects
  \texttt{dataQuality:"visual"} flags so UI problems become first-class inputs
  alongside the numbers.
  \item \textbf{Deterministic guardrails.} Every raw LLM item passes through
  \texttt{validateItems}: enums and page paths are repaired, cited numbers are
  grounded against the real data (ungrounded claims downgraded or dropped when
  \texttt{GROUNDING\_STRICT=true}), and duplicates are folded.
  \texttt{forecastImpact} then attaches a \emph{deterministic}, data-derived
  revenue estimate --- never the LLM's own guess.
  \item \textbf{On-demand code generation.} \texttt{/code/generate} pulls the
  real theme files for the affected page from \texttt{shopify-pp}, caps total
  theme content at 150K chars (OOM defense), and asks a code-strong model
  (Gemini by default) for a \emph{surgical} patch: only the changed lines, plus
  verbatim anchor text so a developer --- or a script --- can locate the exact
  insertion point.
\end{itemize}

\subsection{Service 4 --- Review Hub (\texttt{review-hub})}

\textbf{Stack:} React + Vite, Recharts, \texttt{localStorage} persistence.

\begin{itemize}[leftmargin=1.4em]
  \item \textbf{Two personas.} PO/admin sees the full dashboard and the
  approve/reject controls; developers see only approved suggestions \emph{with}
  their code patches.
  \item \textbf{Resilient by design.} If the backend services are down, the
  dashboard falls back to an embedded demo dataset, so it is always demoable.
  The active data source (\texttt{ai} / \texttt{mock} / \texttt{demo}) is shown
  in the header.
  \item \textbf{Zero-backend state.} Review status persists in
  \texttt{localStorage}; the app talks directly to the Analytics and AI
  services from the browser.
\end{itemize}

% ═════════════════════════════════════════════════════════════════════════════
\section{Product Management \& Delivery}
% ═════════════════════════════════════════════════════════════════════════════

Beyond the engineering, the project was run as a structured product engagement
--- the artifacts below are as much a part of the deliverable as the code.

\subsection{Agile backlog}

The backlog is a three-tier Agile hierarchy (\textbf{Epics $\rightarrow$
Features $\rightarrow$ User Stories}), using \textbf{MoSCoW} prioritization and
\textbf{Given--When--Then} (GWT) acceptance criteria throughout. It is
structured across four epics:

\begin{center}
\begin{tabularx}{\textwidth}{@{}l Y l l@{}}
\toprule
\textbf{Epic} & \textbf{Description} & \textbf{Status} & \textbf{Scope}\\
\midrule
E-001 & Data ingestion \& normalization & Complete & MVP (v1)\\
E-002 & Analytics Service REST API (Node/Express) & Complete & MVP (v1)\\
E-003 & AI Suggestion Service --- LLM recommendation generation & Complete & MVP (v1)\\
E-004 & Code generation / developer-facing patch output & Complete & MVP (v1)\\
\bottomrule
\end{tabularx}
\end{center}

All four epics are complete and in MVP scope. E-004 (developer-facing code
generation) is realized as the AI service's \texttt{/code/generate} endpoint,
which fetches the real theme code for the affected page from the extractor and
generates a surgical, ready-to-paste patch.

\subsection{Key product decisions (and their rationale)}

\begin{itemize}[leftmargin=1.4em]
  \item \textbf{Two LLM roles, deliberately split.} Fast, real-time
  \emph{suggestion} generation vs.\ slower, code-strong \emph{patch}
  generation are different jobs and use different models --- this split should
  not be collapsed into ``just use one model''. (Handover pinned Groq
  Llama\,3.3\,70B for suggestions and Gemini for code; the current code
  generalizes suggestions to a provider-agnostic layer with optional
  multi-model orchestration, and keeps Gemini as the default code model.)
  \item \textbf{Confidence capping over fabrication.} Some Shopify metrics
  (bounce rate, cart-abandonment rate, average load time) can return
  \texttt{null}. Rather than estimate or fabricate a value, the system
  \emph{caps} the confidence of any recommendation touching those metrics. This
  is a principled choice to avoid hallucinated numbers --- framed to merchants
  as a \emph{managed limitation}, not an unresolved gap.
  \item \textbf{Detectors skip silently when data is missing.} The optional
  detectors (\texttt{cta\_below\_fold}, \texttt{poor\_quality\_traffic\_source})
  depend on crawler data that only exists when a website URL is supplied. With
  no URL they skip rather than guess --- no hallucinated output when the
  underlying data is absent.
  \item \textbf{Approval as the hand-off trigger.} A PO approving a
  recommendation is precisely what promotes it into the developer view as
  implementation-ready --- the mechanism that keeps a human between the AI and
  the storefront.
  \item \textbf{Mock data is always labelled.} Demo/unsynced accounts see mock
  data clearly marked as test-store data, so it is never confusable with a real
  merchant's live analytics.
\end{itemize}

\subsection{Competitive \& discovery context}

Product discovery benchmarked the engine against experience-analytics and CRO
tools (\emph{Contentsquare}, \emph{VWO}, \emph{AICRO.io}), AI coding assistants
(\emph{Cursor}) as a reference for the code-generation experience, and
comparable consultancy players (\emph{Accenture}, \emph{Infosys}, \emph{EPAM},
\emph{Thoughtworks}). The differentiator that emerged: not another analytics
dashboard, but the \emph{translation layer} from measurement to a reviewed,
implementable code change.

% ═════════════════════════════════════════════════════════════════════════════
\section{How to Connect All Services}\label{sec:connect}
% ═════════════════════════════════════════════════════════════════════════════

The services are wired together entirely through URLs in environment variables
and one frontend config file. Getting these right is the whole integration.

\subsection{The wiring map}

\begin{center}
\begin{tabularx}{\textwidth}{@{}l l Y@{}}
\toprule
\textbf{Consumer} & \textbf{Setting} & \textbf{Points at}\\
\midrule
\texttt{shopify-pp} & \texttt{ANALYTICS\_SERVICE\_URL} (env) & Analytics \texttt{/api/analytics/ingest}\\
\texttt{ai service} & \texttt{ANALYTICS\_SERVICE\_URL} (env) & Analytics base URL\\
\texttt{ai service} & \texttt{SHOPIFY\_APP\_URL} (env) & Extractor base URL (for theme code)\\
\texttt{review-hub} & \texttt{ANALYTICS\_URL} (config.js) & Analytics base URL\\
\texttt{review-hub} & \texttt{AI\_URL} (config.js) & AI service base URL\\
\texttt{review-hub} & \texttt{SHOPIFY\_APP\_URL} (config.js) & Extractor base URL\\
\bottomrule
\end{tabularx}
\end{center}

\subsection{Local development}

Install dependencies once per service, then start each in its own terminal in
dependency order (Analytics first, since AI depends on it):

\begin{lstlisting}[language=bash]
# 1. Analytics Service  -> http://localhost:4000  (no env vars needed)
cd "analytic service" && npm install && npm start

# 2. AI Service         -> http://localhost:5001
cd "ai service" && npm install
cp .env.example .env      # add GROQ / MISTRAL / CEREBRAS / GEMINI keys
npm start

# 3. Shopify Extractor  -> http://localhost:3000
cd shopify-pp && npm install
cp .env.example .env      # add Shopify client id/secret, ENCRYPTION_KEY
npm start

# 4. Review Hub         -> http://localhost:5173
cd review-hub && npm install && npm run dev
\end{lstlisting}

\subsection{Exposing local services with Cloudflare quick tunnels}

Because the Review Hub runs in the \emph{browser}, it cannot reach a
teammate's \texttt{localhost}. A Cloudflare quick tunnel mints a public HTTPS
URL for a local port with no account or deploy required --- useful to dodge
free-tier rate limits on hosted LLM backends:

\begin{lstlisting}[language=bash]
brew install cloudflared
cloudflared tunnel --url http://localhost:4000   # -> Analytics tunnel URL
cloudflared tunnel --url http://localhost:5001   # -> AI tunnel URL
\end{lstlisting}

Paste the printed \texttt{*.trycloudflare.com} URLs into
\texttt{review-hub/src/lib/config.js} (\texttt{ANALYTICS\_URL}, \texttt{AI\_URL})
and \texttt{ai service/.env} (\texttt{ANALYTICS\_SERVICE\_URL}).

\begin{keybox}{Two gotchas that cost real debugging time}
\textbf{(1) Tunnel URLs are single-use} --- each restart mints a brand-new
random URL; you must re-paste it everywhere. \textbf{(2)} \texttt{dotenv} only
reads \texttt{.env} once at startup, so \textbf{restart the AI service after
editing its \texttt{.env}}. Vite's dev server, by contrast, hot-reloads
\texttt{config.js} on save.
\end{keybox}

% ═════════════════════════════════════════════════════════════════════════════
\section{Deployment --- Vercel, Railway, Cloudflare \& More}\label{sec:deploy}
% ═════════════════════════════════════════════════════════════════════════════

Each service has different runtime needs, so no single platform is optimal for
all four. The table below maps each service to its best-fit host and why.

\begin{center}
\begin{tabularx}{\textwidth}{@{}l l Y@{}}
\toprule
\textbf{Service} & \textbf{Recommended host} & \textbf{Why}\\
\midrule
Review Hub & \textbf{Vercel} / Cloudflare Pages & Static Vite build; global CDN; zero server\\
Shopify Extractor & \textbf{Railway} / Fly.io & Long-lived Node process + persistent SQLite volume + webhooks\\
Analytics Service & \textbf{Railway} / Fly.io (Docker) & Needs Playwright/Chromium; ships its own Dockerfile\\
AI Service & \textbf{Railway} / Fly.io & Long-lived Node process; outbound LLM calls; env secrets\\
\bottomrule
\end{tabularx}
\end{center}

\subsection{Review Hub on Vercel (or Cloudflare Pages)}

The Review Hub is a pure static build --- ideal for a CDN host.

\begin{lstlisting}[language=bash]
# From the review-hub/ directory
npm install -g vercel
vercel               # first run links the project
vercel --prod        # deploys the production build
\end{lstlisting}

\begin{itemize}[leftmargin=1.4em]
  \item \textbf{Build command:} \texttt{npm run build} \quad
  \textbf{Output dir:} \texttt{dist}
  \item Because service URLs currently live in \texttt{src/lib/config.js} (not
  env vars), point them at your \emph{deployed} Analytics/AI URLs before
  building. \emph{Recommended improvement:} read them from \texttt{import.meta.env.VITE\_*}
  so the same build works across environments (see Section~\ref{sec:improve}).
  \item \textbf{Cloudflare Pages} is an equivalent alternative: connect the
  repo, set build command \texttt{npm run build}, output \texttt{dist}.
\end{itemize}

\subsection{Shopify Extractor \& AI Service on Railway}

Railway builds Node apps with Nixpacks automatically and provides persistent
volumes --- essential for the extractor's SQLite database.

\begin{lstlisting}[language=bash]
npm install -g @railway/cli
railway login
railway init                 # run inside each service folder
railway up                   # build & deploy
\end{lstlisting}

\begin{itemize}[leftmargin=1.4em]
  \item \textbf{Start command:} \texttt{npm start} for both.
  \item \textbf{Extractor env vars:} \texttt{SHOPIFY\_CLIENT\_ID},
  \texttt{SHOPIFY\_CLIENT\_SECRET}, \texttt{SHOPIFY\_SCOPES}, \texttt{APP\_URL}
  (the public Railway URL), \texttt{ANALYTICS\_SERVICE\_URL},
  \texttt{ENCRYPTION\_KEY}, and \texttt{DB\_PATH} pointing at a
  \textbf{mounted volume} (e.g.\ \texttt{/data/shops.db}) --- otherwise every
  redeploy wipes installed shops.
  \item \textbf{AI env vars:} the provider API keys, \texttt{AI\_PROVIDER},
  \texttt{ANALYTICS\_SERVICE\_URL}, \texttt{SHOPIFY\_APP\_URL}.
  \item \textbf{Update the Shopify app config:} \texttt{shopify.app.toml}'s
  \texttt{application\_url} and \texttt{redirect\_urls} must match the deployed
  extractor URL.
\end{itemize}

\subsection{Analytics Service via Docker (Railway / Fly.io / any container host)}

The Analytics Service needs Chromium, so it ships a Dockerfile built on the
official Playwright image:

\begin{lstlisting}[language=bash]
FROM mcr.microsoft.com/playwright:v1.61.1-jammy
WORKDIR /app
COPY package*.json ./
RUN npm ci --omit=dev
COPY . .
CMD ["node", "index.js"]
\end{lstlisting}

\begin{lstlisting}[language=bash]
# Railway auto-detects the Dockerfile; or on Fly.io:
fly launch --dockerfile Dockerfile
fly deploy
\end{lstlisting}

\begin{itemize}[leftmargin=1.4em]
  \item Optionally set \texttt{AWS} S3 credentials so screenshots persist in
  object storage instead of the container's ephemeral filesystem.
  \item Ensure the host allocates enough memory --- Chromium is memory-hungry;
  Railway's smallest tier can OOM on image-heavy crawls.
\end{itemize}

\subsection{Cloudflare's role}

Cloudflare serves two distinct purposes here, which are easy to confuse:

\begin{enumerate}[leftmargin=1.4em]
  \item \textbf{Quick tunnels} (\texttt{cloudflared tunnel --url}) --- an
  ephemeral dev bridge from a public URL to a local port
  (Section~\ref{sec:connect}). Not for production.
  \item \textbf{Cloudflare Pages} --- a production static host for the Review
  Hub, equivalent to Vercel.
\end{enumerate}

Note that \textbf{Playwright-based crawling and long-lived Express servers
cannot run on Cloudflare Workers} (no full Node/Chromium runtime), which is why
the backend services target Railway/Fly rather than Cloudflare's edge.

\subsection{Deployment checklist}

\begin{enumerate}[leftmargin=1.4em]
  \item Deploy Analytics (Docker) $\rightarrow$ note its URL.
  \item Deploy AI Service with \texttt{ANALYTICS\_SERVICE\_URL} set $\rightarrow$ note its URL.
  \item Deploy Extractor with \texttt{ANALYTICS\_SERVICE\_URL} + \texttt{APP\_URL} + persistent \texttt{DB\_PATH}.
  \item Update \texttt{shopify.app.toml} URLs to the deployed extractor.
  \item Build \& deploy Review Hub with all three URLs configured.
  \item Smoke-test the pipeline using \texttt{TESTING.md}'s end-to-end order.
\end{enumerate}

% ═════════════════════════════════════════════════════════════════════════════
\section{How Can We Improve It}\label{sec:improve}
% ═════════════════════════════════════════════════════════════════════════════

The current system is a strong, demoable prototype. The improvements below are
grouped by theme and roughly ordered by leverage.

\subsection{Persistence \& state}
\begin{itemize}[leftmargin=1.4em]
  \item \textbf{Replace in-memory caches with a real datastore.} Both the
  Analytics cache and the AI service's report/suggestion state are in-memory
  objects --- a restart wipes everything. Move to Postgres or Redis so reports,
  review decisions, and suggestion status survive redeploys and scale beyond
  one process.
  \item \textbf{Move review state off \texttt{localStorage}.} The Review Hub's
  approve/reject decisions live only in the browser; they should live server-side
  so decisions are shared across users and devices and auditable.
\end{itemize}

\subsection{Configuration \& security}
\begin{itemize}[leftmargin=1.4em]
  \item \textbf{Environment-driven frontend config.} Hard-coded URLs in
  \texttt{config.js} require a rebuild per environment; switch to
  \texttt{VITE\_*} env vars so one artifact deploys everywhere.
  \item \textbf{Real authentication.} The Review Hub uses hard-coded users.
  Replace with proper auth (Shopify session tokens, or an SSO provider) and
  role-based access control.
  \item \textbf{Lock down CORS.} Services currently allow all origins; scope
  them to known frontends in production.
\end{itemize}

\subsection{Reliability \& observability}
\begin{itemize}[leftmargin=1.4em]
  \item \textbf{Async job queue for crawling \& generation.} Crawls and LLM
  calls are slow and rate-limited; a queue (BullMQ / a managed queue) with
  retries and backpressure beats fire-and-forget async.
  \item \textbf{Structured logging, metrics, and tracing.} Today's insight is
  \texttt{console.log}; add request tracing across the four services so a
  failed report can be diagnosed end-to-end.
  \item \textbf{Cost \& quota telemetry.} Token usage is already tracked in
  \texttt{meta}; surface it as a dashboard so provider costs and rate-limit
  headroom are visible before they cause 503s.
\end{itemize}

\subsection{AI quality}
\begin{itemize}[leftmargin=1.4em]
  \item \textbf{Close the feedback loop.} Feed PO approve/reject decisions back
  as training signal (few-shot examples or fine-tuning) so suggestion quality
  improves over time per merchant.
  \item \textbf{Patch verification before it reaches a developer.} Lint or
  dry-apply the generated patch against the fetched theme files, so a developer
  never receives a patch whose anchor text no longer matches.
  \item \textbf{Regression suite for detectors and prompts.} The fixtures
  directory is a good start; grow it into a golden-dataset test so prompt or
  threshold changes are measured, not guessed.
\end{itemize}

% ═════════════════════════════════════════════════════════════════════════════
\section{Future Scope}\label{sec:future}
% ═════════════════════════════════════════════════════════════════════════════

\subsection{Immediate next steps (open items)}

With MVP scope complete, the open items are next-phase work rather than
unfinished MVP tasks. In priority order:

\begin{enumerate}[leftmargin=1.4em]
  \item \textbf{Run the first merchant pilot.} The single most important next
  step --- v1 is only validated once one recommendation is implemented on a
  real store and its impact measured. Everything else is subordinate to this.
  \item \textbf{Harden code generation from the pilot.} E-004 is implemented;
  the pilot is what proves the PO-approval$\rightarrow$patch flow on real theme
  code and surfaces the rough edges.
  \item \textbf{Add data depth (GA4 and beyond)} to reduce reliance on
  confidence capping --- the first concrete step toward the universal-connector
  vision below.
  \item \textbf{Specify rate limits concretely.} ``Generous limits'' was
  flagged during backlog review as under-specified; real numbers are needed
  before any external-facing rollout.
  \item \textbf{Frame crawler optionality correctly.} When explaining why the
  two optional detectors need a website URL, frame it as a barrier-to-entry and
  a time-vs-quality tradeoff for the merchant --- not around token limits.
\end{enumerate}

\subsection{Vision}

Today the pipeline is Shopify-specific and ends by handing a developer a patch
to copy-paste. The future vision inverts both limitations:

\begin{keybox}{Where this is going}
A library of lightweight \textbf{extensions / connectors} pulls \emph{raw data}
from any analytics platform. That data is turned into ranked suggestions and
generated code \emph{exactly as today}. But instead of a human copy-pasting the
patch, the approved change is \textbf{applied to the codebase automatically} ---
committed, opened as a pull request, and merged --- with no manual step. And it
is \textbf{no longer limited to Shopify.}
\end{keybox}

\subsection{Pillar 1 --- Universal connectors (beyond Shopify)}

The extractor becomes one of many pluggable \emph{connectors} behind a common
interface. Each connector's only job: authenticate to a platform and emit raw
data in the shape the Analytics Service already expects
(\texttt{overview}, \texttt{pages}, \texttt{trafficSources}, \texttt{devices},
\texttt{funnel}, Core Web Vitals). Because the ingest contract is already
platform-neutral, everything downstream --- detectors, AI, review --- works
unchanged. Candidate connectors:

\begin{center}
\begin{tabularx}{\textwidth}{@{}l Y@{}}
\toprule
\textbf{Connector} & \textbf{Raw data it would emit}\\
\midrule
Google Analytics 4 & Sessions, conversions, funnels, traffic sources, devices\\
Adobe Analytics & Enterprise session \& funnel data\\
WooCommerce / Magento / BigCommerce & Non-Shopify e-commerce metrics + theme/template code\\
Plausible / Matomo & Privacy-first web analytics\\
Mixpanel / Amplitude & Product-analytics event funnels\\
Search Console / PageSpeed & SEO \& Core Web Vitals at scale\\
\bottomrule
\end{tabularx}
\end{center}

The connector interface is deliberately thin: \texttt{authenticate()},
\texttt{extract()} $\rightarrow$ normalized payload, and optionally
\texttt{getSourceCode(page)} for platforms where the code is fetchable. A
connector that cannot supply source code simply skips the auto-apply pillar and
still produces suggestions.

\subsection{Pillar 2 --- Automatic code application}

The single biggest manual step today is a developer copy-pasting the patch. The
patch format was \emph{designed} to remove that step: each patch already carries
verbatim anchor text (``insert after this exact line \dots'') precisely so a
\emph{script} --- not just a human --- can locate the edit. The roadmap:

\begin{enumerate}[leftmargin=1.4em]
  \item \textbf{Anchor-based patch applier.} A deterministic engine locates the
  anchor by exact string match and applies the change, failing loudly if the
  anchor drifted (never a fuzzy guess on live code).
  \item \textbf{Git integration.} For code-backed platforms, the applier
  commits to a branch and opens a \textbf{pull request} --- the PO's ``approve''
  becomes the trigger, and the human review gate moves to standard PR review.
  \item \textbf{Shopify theme write-back.} Using the Shopify Theme API, apply
  the change to a \emph{duplicated draft theme} first, screenshot the result,
  and only publish after approval.
  \item \textbf{Verify-before-merge.} Re-run the crawler on the changed page
  and confirm the flagged problem is actually resolved before the PR is allowed
  to merge --- closing the observe$\rightarrow$act$\rightarrow$verify loop
  automatically.
\end{enumerate}

\subsection{Pillar 3 --- Continuous optimization}

With connectors pulling data on a schedule and patches applying themselves, the
system becomes a \emph{continuous} optimization engine rather than a one-shot
report: detect $\rightarrow$ suggest $\rightarrow$ (PO-approve) $\rightarrow$
apply $\rightarrow$ measure $\rightarrow$ repeat. A/B testing each applied change
against a holdout would let the system \emph{prove} lift rather than forecast
it, and feed real outcomes back into suggestion ranking.

\subsection{Target architecture}

\begin{center}
\begin{tabularx}{\textwidth}{@{}l Y@{}}
\toprule
\textbf{Layer} & \textbf{Evolution}\\
\midrule
Connectors & Shopify extractor $\rightarrow$ pluggable connector SDK (GA4, Woo, Adobe, \dots)\\
Analytics & Unchanged contract; add scheduled/continuous ingestion\\
AI & Unchanged generation; add feedback-trained ranking\\
Apply & \emph{New}: anchor-based applier + Git/Theme write-back + verify\\
Review & PO approval becomes the PR/apply trigger, not a copy-paste handoff\\
\bottomrule
\end{tabularx}
\end{center}

\begin{keybox}{The through-line}
Every future pillar reuses the platform-neutral ingest contract and the
anchor-carrying patch format that already exist. The prototype was, in effect,
built as the Shopify \emph{instance} of a general system --- the roadmap is
mostly about widening the front door (connectors) and automating the last step
(apply), not rewriting the middle.
\end{keybox}

% ═════════════════════════════════════════════════════════════════════════════
\section{Conclusion}
% ═════════════════════════════════════════════════════════════════════════════

Exp Intelligence demonstrates that the valuable, hard part of storefront
optimization is not measurement but the \emph{translation from measurement to
implementation}. By chaining statistical rigor, multimodal (numbers + visual)
reasoning, and theme-grounded code generation --- all behind a human review
gate --- it turns a wall of analytics into a short, ranked list of actions with
the code to ship them. The architecture's platform-neutral ingest contract and
script-locatable patch format mean the ambitious future --- universal
connectors and automatic code application beyond Shopify --- is an extension of
what exists, not a rebuild.

\vfill
\begin{center}
{\color{rulegray}\rule{0.5\textwidth}{0.5pt}}\\[0.3cm]
{\small\itshape End of report}
\end{center}

\end{document}
